\documentclass[tikz,border=10pt]{standalone}
\usepackage{amsmath}
\usepackage[x11names]{xcolor}
\definecolor{bgdark}{HTML}{171D2B}
\definecolor{aksen}{HTML}{F7B538}
\begin{document}
\pagecolor{bgdark}
\color{white}
\begin{tikzpicture}[scale=1.4, line cap=round]
  % O=(0,0), r=2, P=(4,3): d=5, alpha=36.87deg, beta=asin(2/5)=23.58deg
  % T1 pada sudut 60.45deg, T2 pada sudut 13.29deg
  \coordinate (O) at (0,0);
  \coordinate (P) at (4,3);
  \coordinate (T1) at (0.986,1.741);
  \coordinate (T2) at (1.946,0.460);
  % lingkaran
  \draw[thick, white] (O) circle (2);
  % garis singgung (aksen kuning)
  \draw[aksen, very thick] (P) -- (T1);
  \draw[aksen, very thick] (P) -- (T2);
  % jari-jari ke titik singgung + garis OP putus-putus
  \draw[white] (O) -- (T1);
  \draw[dashed, gray!60] (O) -- (P);
  \pic [draw, white, angle radius=7pt] {right angle = O--T1--P};
  % titik dan label
  \fill (O) circle (1.6pt) node[below left] {$O$};
  \fill (P) circle (1.6pt) node[above right] {$P$};
  \fill (T1) circle (1.6pt) node[above left] {$T_1$};
  \fill (T2) circle (1.6pt) node[below right] {$T_2$};
  \node[gray!50, left=2pt] at ($(O)!0.5!(T1)$) {$r$};
\end{tikzpicture}
\end{document}
